% ============================== CHAPTER 5 ============================
\clearpage
\setlecture{Lecture V}\renewcommand{\tabname}{V}
\chapter{Applications: the Derivative as a Linear Map}\label{ch:linear}

So far our functions returned a single number. Many important objects return several
numbers at once: a change of coordinates $(r,\theta)\mapsto(x,y)$, the position of a robot's
hand as a function of its joint angles, a velocity field. This chapter defines the derivative
of such a mapping $F:\R^n\to\R^m$. The answer is beautiful: the derivative is a
\emph{linear map}, represented by the \emph{Jacobian matrix}. With it, the chain rule becomes a
matrix product, coordinate changes become transparent, and we can prove two of the deepest
results of the course: the inverse mapping theorem and the implicit function theorem.

\section{The derivative as a linear map}

\subsection{From differentiability to a linear map}

Let $U\subset\R^n$ be open, $P\in U$, and $f$ differentiable at $P$. For small $H$,
\[ f(P+H)=f(P)+A\cdot H+\norm Hg(H),\qquad \lim_{\norm H\to0}g(H)=0, \]
where $A=\grad f(P)=\big(D_1f(P),\dots,D_nf(P)\big)$. The map $L=L_A$ defined by
\[ L(H)=A\cdot H \]
is \emph{linear}: $L(H+K)=L(H)+L(K)$ and $L(cH)=cL(H)$. So differentiability reads
\[ f(P+H)=f(P)+L(H)+\norm Hg(H). \]

\begin{intuition}
In one variable, $f(a+h)=f(a)+ch+\abs h g(h)$ with $c=f'(a)$, and ``multiplication by $c$'' is a
linear map $L_c(h)=ch$. The derivative was always secretly a linear map; for functions of one
variable we just identified the map $h\mapsto ch$ with the number $c$. In several variables this
identification is no longer convenient, so we keep the map itself.
\end{intuition}

\begin{definition}{Derivative of a scalar function}
The \emph{derivative} of $f$ at $P$ is the linear map $H\mapsto\grad f(P)\cdot H$. It is
denoted $Df(P)$ or $f'(P)$:
\[ f(P+H)=f(P)+Df(P)H+\norm Hg(H), \]
or equivalently
\[ f(P+H)=f(P)+f'(P)H+\norm Hg(H). \]
\end{definition}

\subsection{Differentiable mappings}

\begin{definition}{Differentiability of a mapping}
Let $U$ be open in $\R^n$, $F:U\to\R^m$ a mapping and $P\in U$. $F$ is \emph{differentiable at
$P$} if there exist a linear map $L:\R^n\to\R^m$ and a mapping $G$, defined for all
sufficiently small vectors $H$, such that
\[ F(P+H)=F(P)+LH+\norm H\,G(H)\qquad\text{and}\qquad \lim_{\norm H\to0}G(H)=O . \]
\end{definition}

\begin{definition}{Tangent map and little-o}
A linear map $L$ as above is said to be \emph{tangent} to $F$ at $P$: it is the best linear
approximation of $F$ near $P$. A map $\psi$ defined for small $H$ is $o(H)$ (``little oh of
$H$'') if
\[ \lim_{\norm H\to0}\frac{\psi(H)}{\norm H}=O . \]
Thus differentiability says $F(P+H)=F(P)+L(H)+o(H)$.
\end{definition}

\subsection{Uniqueness}

\begin{theorem}[label=thm:unique]{Uniqueness of the tangent linear map}
If two linear maps $L$ and $M$ are both tangent to $F$ at $P$, then $L=M$.
\end{theorem}

\begin{intuition}
Two linear maps that agree up to an error much smaller than $\norm H$ on \emph{all} small vectors
must agree exactly. Scaling makes this precise: a linear map scales with $t$, while the error
becomes negligible as $t\to0$.
\end{intuition}

\begin{proof}
Write both expansions:
\begin{align*} F(P+H)&=F(P)+LH+\norm HG_1(H),\\ F(P+H)&=F(P)+MH+\norm HG_2(H), \end{align*}
with $G_1,G_2\to O$. Fix any vector $Y\neq O$ and evaluate both at $H=tY$, $t>0$ small. Subtracting,
\[ O=L(tY)-M(tY)+\norm{tY}\big[G_1(tY)-G_2(tY)\big]. \]
By linearity $L(tY)=tL(Y)$, $M(tY)=tM(Y)$, and $\norm{tY}=t\norm Y$. Dividing by $t$ and putting
$G=G_1-G_2$,
\[ M(Y)-L(Y)=\norm Y\,G(tY). \]
The left side does not depend on $t$. Let $t\to0$: then $tY\to O$, so $G(tY)\to O$, and the right
side tends to $O$. Hence $M(Y)=L(Y)$. This holds for every $Y$ (trivially for $Y=O$), so $M=L$.
\end{proof}

By uniqueness, \emph{the} tangent map is called the \textbf{derivative} of $F$ at $P$ and denoted
$F'(P)$ or $DF(P)$:
\[ F(P+H)=F(P)+F'(P)H+\norm H\,G(H),\qquad G(H)\to O . \]

\begin{keyidea}
The derivative of $F:\R^n\to\R^m$ at $P$ is a \textbf{linear map} $\R^n\to\R^m$: it eats a small
displacement $H$ of the input and returns the corresponding (first-order) displacement of the
output. It is not merely a list of partial derivatives.
\end{keyidea}

\begin{example}{linear and affine maps}
If $F(X)=AX+B$ with $A$ an $m\times n$ matrix, then $F(P+H)=F(P)+AH$ exactly, so the remainder
is $0$ and $F'(P)=A$ for every $P$. A linear map is its own derivative, as a straight line is its
own tangent.
\end{example}

\section{The Jacobian matrix}

\subsection{Coordinates}

Let $F:\R^n\to\R^m$ have coordinate functions $f_1,\dots,f_m$:
\[ F(X)=\begin{pmatrix}f_1(X)\\ \vdots\\ f_m(X)\end{pmatrix}=\tr\big(f_1(X),\dots,f_m(X)\big), \]
where we view $X=\tr(x_1,\dots,x_n)$ as a column vector, and assume the partial derivatives of
each $f_i$ exist.

\begin{definition}{Jacobian matrix}
The \emph{Jacobian matrix} of $F$ at $X$ is the $m\times n$ matrix of partial derivatives
\[
J_F(X)=\left(\pd{f_i}{x_j}\right)=\begin{pmatrix}D_1f_1(X)&\cdots&D_nf_1(X)\\ \vdots&&\vdots\\ D_1f_m(X)&\cdots&D_nf_m(X)\end{pmatrix}.
\]
Row $i$ is the gradient of $f_i$; column $j$ records how all outputs react to $x_j$. For
$F(x,y)=(f(x,y),g(x,y))$,
\[ J_F(x,y)=\begin{pmatrix}\pd fx&\pd fy\\[3pt] \pd gx&\pd gy\end{pmatrix}. \]
\end{definition}

\begin{center}
\begin{tikzpicture}
  \matrix (m) [matrix of math nodes, left delimiter=(, right delimiter=), nodes={minimum width=1.25cm, minimum height=0.7cm, font=\small}]
  {  \partial_1 f_1 & \partial_2 f_1 & \cdots & \partial_n f_1 \\
     \partial_1 f_2 & \partial_2 f_2 & \cdots & \partial_n f_2 \\
     \vdots & \vdots &  & \vdots \\
     \partial_1 f_m & \partial_2 f_m & \cdots & \partial_n f_m \\};
  \begin{scope}[on background layer]
   \fill[coral, opacity=0.18] (m-2-1.north west) rectangle (m-2-4.south east);
   \fill[cyanbp, opacity=0.15] (m-1-2.north west) rectangle (m-4-2.south east);
  \end{scope}
  \node[font=\scriptsize, text=coraldark, anchor=west] at ([xshift=0.5cm]m-2-4.east) {row $i$ $=\nabla f_i$};
  \node[font=\scriptsize, text=cyanbp, anchor=north] at ([yshift=-0.2cm]m-4-2.south) {column $j$: effect of $x_j$};
\end{tikzpicture}
\end{center}

\begin{example}{a map $\R^2\to\R^2$}
Let $F(x,y)=\begin{pmatrix}x^2+y^2\\ e^{xy}\end{pmatrix}$. Then
\[ J_F(x,y)=\begin{pmatrix}2x&2y\\ ye^{xy}&xe^{xy}\end{pmatrix},\qquad J_F(1,1)=\begin{pmatrix}2&2\\ e&e\end{pmatrix}. \]
\end{example}

\begin{example}{a map $\R^2\to\R^3$}
For $F(x,y)=\begin{pmatrix}xy\\ \sin x\\ x^2y\end{pmatrix}$,
\[ J_F(x,y)=\begin{pmatrix}y&x\\ \cos x&0\\ 2xy&x^2\end{pmatrix},\qquad
J_F\!\left(\pi,\tfrac\pi2\right)=\begin{pmatrix}\pi/2&\pi\\ -1&0\\ \pi^2&\pi^2\end{pmatrix}. \]
A $3\times2$ matrix: three outputs, two inputs.
\end{example}

\subsection{Derivative and Jacobian}

\begin{theorem}[label=thm:jac]{The derivative is the Jacobian}
Let $U$ be open in $\R^n$ and $F:U\to\R^m$ with coordinate functions $f_1,\dots,f_m$. If each
$f_i$ is differentiable at $X\in U$, then $F$ is differentiable at $X$ and the matrix
representing $F'(X)=DF(X)$ is the Jacobian matrix $J_F(X)$:
\[ F(X+H)=F(X)+J_F(X)H+o(H). \]
\end{theorem}

\begin{proof}
\textbf{Step 1.} For each coordinate function,
\[ f_i(X+H)=f_i(X)+\grad f_i(X)\cdot H+\norm Hg_i(H),\qquad g_i(H)\to0 . \]
\textbf{Step 2: stack the $m$ equations into a column.} The terms $\grad f_i(X)\cdot H$ are the
entries of the product $J_F(X)H$, since row $i$ of $J_F(X)$ is $\grad f_i(X)$. With
$G(H)=\tr\big(g_1(H),\dots,g_m(H)\big)$,
\[ F(X+H)=F(X)+J_F(X)H+\norm H\,G(H). \]
\textbf{Step 3.} Each component of $G(H)$ tends to $0$, so $G(H)\to O$. Thus the linear map
$H\mapsto J_F(X)H$ is tangent to $F$ at $X$, and by uniqueness it is $F'(X)$.
\end{proof}

\section{Differentials and coordinate changes}

\subsection{The differential of a scalar function}

\begin{definition}{Differential}
Let $f:\R^n\to\R$ be differentiable at $P$. The \emph{differential} of $f$ at $P$ is the derivative
viewed as a linear functional,
\[ df_P:=Df(P):\R^n\to\R,\qquad df_P(H)=\sum_{i=1}^n\pd f{x_i}(P)\,h_i . \]
\end{definition}

\begin{trap}
$df_P$ is \textbf{not} an ``infinitesimally small number''. It is a linear map that acts on an
increment $H$ and returns a number. ``$df\approx\Delta f$'' means $df_P(H)\approx f(P+H)-f(P)$ for small $H$.
\end{trap}

For each coordinate function $x_i(X)=x_i$, the derivative is the linear functional
$dx_i(H)=h_i$: the differentials $dx_1,\dots,dx_n$ simply \emph{extract the components} of $H$.

\begin{proposition}{Coordinate expression of the differential}
For a differentiable scalar function $f$,
\[ df_P=\sum_{i=1}^n\pd f{x_i}(P)\,dx_i\qquad\text{as an identity of linear functionals.} \]
In two variables: $df=f_x\,dx+f_y\,dy$.
\end{proposition}

\begin{proof}
Apply both sides to $H$: the left gives $\sum\pd f{x_i}(P)h_i$; the right gives
$\sum\pd f{x_i}(P)\,dx_i(H)=\sum\pd f{x_i}(P)h_i$. Equal for every $H$, so equal as maps.
\end{proof}

\begin{example}{error estimate with differentials}
The volume of a cylinder is $V=\pi r^2h$, so $dV=2\pi rh\,dr+\pi r^2\,dh$. For $r=5$ cm,
$h=10$ cm measured with errors $\abs{dr},\abs{dh}\le0.1$ cm:
\[ \abs{dV}\le2\pi(5)(10)(0.1)+\pi(25)(0.1)=10\pi+2.5\pi=12.5\pi\approx39\ \text{cm}^3, \]
out of $V=250\pi\approx785$ cm$^3$: about $5\%$. Note that the error in $r$ costs four times more
than the error in $h$.
\end{example}

\subsection{Differentials under a change of coordinates}

Let $F(u_1,\dots,u_n)=\big(x_1(u),\dots,x_m(u)\big)$. Since $J_F(u)$ represents $DF(u)$,
\[
dx_i=\sum_{j=1}^n\pd{x_i}{u_j}\,du_j,\qquad\text{equivalently}\qquad
\begin{pmatrix}dx_1\\ \vdots\\ dx_m\end{pmatrix}=J_F(u)\begin{pmatrix}du_1\\ \vdots\\ du_n\end{pmatrix}.
\]
The Jacobian matrix tells us how coordinate differentials transform locally.

\begin{example}{polar-coordinate differentials}
For $x=r\cos\theta$, $y=r\sin\theta$:
\[ dx=\cos\theta\,dr-r\sin\theta\,d\theta,\qquad dy=\sin\theta\,dr+r\cos\theta\,d\theta, \]
that is
\[
\begin{pmatrix}dx\\dy\end{pmatrix}=\begin{pmatrix}\cos\theta&-r\sin\theta\\ \sin\theta&r\cos\theta\end{pmatrix}\begin{pmatrix}dr\\d\theta\end{pmatrix},
\]
and $\det J_F(r,\theta)=r\cos^2\theta+r\sin^2\theta=r$.
\end{example}

\subsection{Area and volume elements}

A linear map multiplies areas (or volumes) by the absolute value of its determinant. For a
differentiable map $x=F(u)$, a tiny box in the $u$-coordinates is mapped to an almost-box (almost a
parallelogram) whose volume is multiplied by $\abs{\det J_F(u)}$. Hence, locally,
\[ \boxed{\,dV_x=\abs{\det J_F(u)}\,dV_u\,} \]
and in two dimensions $dA_x=\abs{\det J_F(u,v)}\,dA_{u,v}$.

\begin{center}
\begin{tikzpicture}[scale=1]
  \begin{scope}
    \draw[bpgrid, step=0.5] (-0.3,-0.3) grid (3.3,2.8);
    \draw[navy, ->, line width=0.6pt] (-0.3,0) -- (3.3,0) node[right, font=\small]{$r$};
    \draw[navy, ->, line width=0.6pt] (0,-0.3) -- (0,2.8) node[above, font=\small]{$\theta$};
    \fill[coral, opacity=0.35] (1.8,1.2) rectangle (2.3,1.6);
    \draw[coral, line width=1pt] (1.8,1.2) rectangle (2.3,1.6);
    \node[font=\scriptsize, text=coraldark, below] at (2.05,1.2) {$dr$};
    \node[font=\scriptsize, text=coraldark, right] at (2.3,1.4) {$d\theta$};
  \end{scope}
  \draw[-{Stealth[length=3mm]}, plum, line width=1.4pt] (3.7,1.3) to[bend left=20] node[above, font=\small]{$F$} (5.0,1.3);
  \begin{scope}[xshift=7.1cm, yshift=0.2cm]
    \draw[bpgrid, step=0.5] (-1.8,-0.5) grid (2.6,2.6);
    \draw[navy, ->, line width=0.6pt] (-1.8,0) -- (2.6,0) node[right, font=\small]{$x$};
    \draw[navy, ->, line width=0.6pt] (0,-0.5) -- (0,2.6) node[above, font=\small]{$y$};
    \fill[coral, opacity=0.35] (40:1.8) arc (40:58:1.8) -- (58:2.3) arc (58:40:2.3) -- cycle;
    \draw[coral, line width=1pt] (40:1.8) arc (40:58:1.8) -- (58:2.3) arc (58:40:2.3) -- cycle;
    \draw[navy, dashed, line width=0.5pt] (0,0) -- (40:2.3);
    \draw[navy, dashed, line width=0.5pt] (0,0) -- (58:2.3);
    \node[font=\scriptsize, text=coraldark] at (49:2.75) {$r\,d\theta$};
    \node[font=\scriptsize, text=coraldark, anchor=north west] at (40:2.05) {$dr$};
  \end{scope}
\end{tikzpicture}

{\small\color{navy} A small rectangle $dr\times d\theta$ becomes a small polar ``rectangle'' with sides
$dr$ and $r\,d\theta$, so $dA=r\,dr\,d\theta$, exactly $\abs{\det J_F}=r$.}
\end{center}

\begin{extra}
One often writes $dx_1\cdots dx_n=\abs{\det J_F(u)}\,du_1\cdots du_n$, or even $dx=\abs{\det J_F(u)}\,du$.
Here ``$dx$'' denotes the \emph{volume element}, not a single coordinate differential. This is
exactly the change-of-variables formula you will use for multiple integrals:
$\iint f\,dx\,dy=\iint f(r\cos\theta,r\sin\theta)\,r\,dr\,d\theta$.
\end{extra}

\begin{definition}{Jacobian determinant}
If $F:U\to\R^n$ maps an open set of $\R^n$ into the \emph{same} dimension, $J_F(X)$ is square and
its determinant is the \emph{Jacobian determinant} of $F$ at $X$, denoted $\Delta_F(X)$.
\end{definition}

\begin{example}{a Jacobian determinant}
For $F(x,y)=(x^2+y^2,\;e^{xy})$,
\[
\Delta_F(x,y)=\begin{vmatrix}2x&2y\\ ye^{xy}&xe^{xy}\end{vmatrix}=2x^2e^{xy}-2y^2e^{xy}=2(x^2-y^2)e^{xy}.
\]
So $\Delta_F(1,1)=0$ and $\Delta_F(1,2)=2(1-4)e^2=-6e^2$.
\end{example}

\begin{trap}
In the lecture slides the entries of the bottom row are printed as $ye^x$ and $xe^y$, which leads
to the value $2e^2-8e$ at $(1,2)$. With $F_2=e^{xy}$, both entries carry $e^{xy}$, and the value is
$-6e^2$. If you meet this on an exam, show the Jacobian matrix explicitly so the grader can follow.
\end{trap}

\begin{example}{the polar coordinate map}
$F(r,\theta)=(r\cos\theta,\;r\sin\theta)$ is defined on all of $\R^2$, though for polar coordinates
we take $r\ge0$. It maps a rectangle $[r_1,r_2]\times[\theta_1,\theta_2]$ onto a circular sector (an
annular piece), as in the picture above, and $\Delta_F(r,\theta)=r$.
\end{example}

\begin{reallife}{a robot arm}
A planar arm has two links of lengths $L_1,L_2$ and joint angles $\theta_1,\theta_2$. The hand is at
\[ F(\theta_1,\theta_2)=\big(L_1\cos\theta_1+L_2\cos(\theta_1+\theta_2),\;L_1\sin\theta_1+L_2\sin(\theta_1+\theta_2)\big). \]
The Jacobian maps joint speeds to hand velocity, and a short computation gives
$\Delta_F=L_1L_2\sin\theta_2$. When $\theta_2=0$ (arm fully stretched) the determinant vanishes: the hand
cannot move along the arm's direction, whatever the motors do. Robotics engineers call this a
\emph{singular configuration} and design controllers to avoid it.
\end{reallife}

\begin{center}
\begin{tikzpicture}[scale=1]
  \draw[bpgrid, step=0.5] (-0.5,-0.4) grid (4.5,3.1);
  \fill[navy] (-0.15,-0.15) rectangle (0.15,0);
  \coordinate (O) at (0,0);
  \coordinate (E) at (35:2.2);
  \coordinate (H) at ($(E)+(85:1.6)$);
  \draw[plum, line width=3pt, line cap=round] (O) -- (E);
  \draw[coral, line width=3pt, line cap=round] (E) -- (H);
  \node[pdot] at (O) {}; \node[pdot] at (E) {}; \node[dot] at (H) {};
  \draw[navy, line width=0.5pt] (0.7,0) arc (0:35:0.7); \node[font=\scriptsize] at (17:0.95) {$\theta_1$};
  \draw[navy, dashed, line width=0.5pt] (E) -- ++(35:0.9);
  \draw[navy, line width=0.5pt] ($(E)+(35:0.55)$) arc (35:85:0.55); \node[font=\scriptsize] at ($(E)+(60:0.8)$) {$\theta_2$};
  \node[font=\scriptsize, text=plum, anchor=north west] at ($(O)!0.5!(E)$) {$L_1$};
  \node[font=\scriptsize, text=coraldark, anchor=west] at ($(E)!0.5!(H)$) {$L_2$};
  \node[font=\scriptsize, text=navy, anchor=south] at (H) {hand $F(\theta_1,\theta_2)$};
\end{tikzpicture}
\end{center}

\section{The chain rule for mappings}

Let $U\subset\R^n$ and $V\subset\R^m$ be open, $F:U\to\R^m$ with $F(U)\subset V$, and
$G:V\to\R^s$. If $F$ is differentiable at $X$ and $G$ at $F(X)$, the derivatives are linear maps
\[ F'(X):\R^n\to\R^m,\qquad G'(F(X)):\R^m\to\R^s . \]

\begin{center}
\begin{tikzpicture}
  \node[draw=plum, line width=1pt, fill=plumtint, ellipse, minimum width=2.2cm, minimum height=1.4cm] (U) at (0,0) {};
  \node[font=\small, text=plum] at (0,-1.45) {$U\subset\R^n$};
  \node[pdot, label={[font=\scriptsize]above:$X$}] (x) at (0,0) {};
  \node[draw=plum, line width=1pt, fill=plumtint, ellipse, minimum width=2.2cm, minimum height=1.4cm] (V) at (4,0) {};
  \node[font=\small, text=plum] at (4,1.1) {$V\subset\R^m$};
  \node[pdot, label={[font=\scriptsize]above:$F(X)$}] (y) at (4,0) {};
  \node[draw=coral, line width=1pt, fill=coraltint, ellipse, minimum width=2.2cm, minimum height=1.4cm] (W) at (8,0) {};
  \node[font=\small, text=coraldark] at (8,-1.45) {$\R^s$};
  \node[dot, label={[font=\scriptsize]above:$G(F(X))$}] (z) at (8,0) {};
  \draw[-{Stealth[length=2.5mm]}, plum, line width=1.1pt] (x) to[bend left=25] node[above, font=\small]{$F$} (y);
  \draw[-{Stealth[length=2.5mm]}, plum, line width=1.1pt] (y) to[bend left=25] node[above, font=\small]{$G$} (z);
  \draw[-{Stealth[length=2.5mm]}, coral, line width=1.1pt, dashed] (x) to[bend right=32] node[below, font=\small, fill=white, inner sep=1pt]{$G\circ F$} (z);
\end{tikzpicture}
\end{center}

\begin{theorem}[label=thm:chainmap]{Chain rule}
Under these hypotheses, $G\circ F$ is differentiable at $X$ and
\[ \boxed{\,(G\circ F)'(X)=G'(F(X))\circ F'(X)\,} \]
In matrices: $J_{G\circ F}(X)=J_G(F(X))\,J_F(X)$, a product of an $s\times m$ and an $m\times n$ matrix.
\end{theorem}

\begin{intuition}
Each map is approximately linear near the relevant point. Composing two approximately linear
maps gives an approximately linear map, and its linear part is the composition of the two linear
parts. The work is to check that the two error terms remain small after composition.
\end{intuition}

\begin{proof}
\textbf{Step 1: the two expansions.} Write $Y=F(X)$. By differentiability,
\begin{align*} F(X+H)&=F(X)+F'(X)H+\norm H\Phi_1(H),\\ G(Y+K)&=G(Y)+G'(Y)K+\norm K\Phi_2(K), \end{align*}
with $\Phi_1(H)\to O$ and $\Phi_2(K)\to O$ (set $\Phi_2(O)=O$).

\textbf{Step 2: choose $K$.} Put $K=F(X+H)-F(X)=F'(X)H+\norm H\Phi_1(H)$. Then
$G(F(X+H))=G(Y+K)$, and
\[ (G\circ F)(X+H)=(G\circ F)(X)+G'(Y)F'(X)H+R(H), \]
using linearity of $G'(Y)$, where $R(H)=\norm H\,G'(Y)\Phi_1(H)+\norm K\Phi_2(K)$.

\textbf{Step 3: the remainder is $o(H)$.} A linear map is bounded: there is $C$ with
$\norm{G'(Y)Z}\le C\norm Z$ and $c$ with $\norm{F'(X)H}\le c\norm H$. So
$\norm{G'(Y)\Phi_1(H)}\le C\norm{\Phi_1(H)}\to0$. Also $\norm K\le(c+\norm{\Phi_1(H)})\norm H$, which is at most
$(c+1)\norm H$ for small $H$; hence $K\to O$ and $\norm K\norm{\Phi_2(K)}/\norm H\le(c+1)\norm{\Phi_2(K)}\to0$.
Therefore $R(H)/\norm H\to O$.

\textbf{Step 4.} So $G'(F(X))F'(X)$ is tangent to $G\circ F$ at $X$, and by uniqueness it is
$(G\circ F)'(X)$.
\end{proof}

\begin{example}{a matrix chain rule}
Let $F(x,y)=(x+y,\;xy)$ and $G(u,v)=(u^2,\;u+v)$. Find $J_{G\circ F}(1,2)$.

\textbf{Step 1.} $F(1,2)=(3,2)$.

\textbf{Step 2.} $J_F=\begin{pmatrix}1&1\\y&x\end{pmatrix}$, so $J_F(1,2)=\begin{pmatrix}1&1\\2&1\end{pmatrix}$;
$J_G=\begin{pmatrix}2u&0\\1&1\end{pmatrix}$, so $J_G(3,2)=\begin{pmatrix}6&0\\1&1\end{pmatrix}$.

\textbf{Step 3: multiply (outer map on the left).}
\[ J_{G\circ F}(1,2)=\begin{pmatrix}6&0\\1&1\end{pmatrix}\begin{pmatrix}1&1\\2&1\end{pmatrix}=\begin{pmatrix}6&6\\3&2\end{pmatrix}. \]
\textbf{Check.} $G\circ F=\big((x+y)^2,\;x+y+xy\big)$ has Jacobian
$\begin{pmatrix}2(x+y)&2(x+y)\\1+y&1+x\end{pmatrix}$, which at $(1,2)$ is $\begin{pmatrix}6&6\\3&2\end{pmatrix}$.
\end{example}

\begin{keyidea}
The chain rule of Chapter~\ref{ch:chain} is the special case $n=1$, $s=1$: for a curve $X(t)$ and a
scalar $f$, $J_{f\circ X}=J_f\,J_X=\grad f(X(t))\cdot X'(t)$, a $1\times m$ row times an $m\times1$ column.
\end{keyidea}

\section{Inverse mappings and implicit functions}

\subsection{$C^1$ maps and invertibility}

\begin{definition}{$C^1$ map and $C^1$-invertibility}
Let $U$ be open in $\R^n$ and $F:U\to\R^n$, $F=(f_1,\dots,f_n)$. If all partial derivatives of all
$f_i$ exist and are continuous, $F$ is a \emph{$C^1$-map}. $F$ is \emph{$C^1$-invertible on $U$} if
$V=F(U)$ is open and there is a $C^1$-map $G:V\to U$ with $G\circ F=\mathrm{id}_U$ and
$F\circ G=\mathrm{id}_V$.
\end{definition}

\begin{example}{translations}
For a fixed vector $A$, $F(X)=X+A$ is $C^1$-invertible on $\R^n$ with inverse $G(Y)=Y-A$.
\end{example}

\begin{example}{polar coordinates on a half-plane}
Let $U=\{(r,\theta):r>0,\ 0<\theta<\pi\}$ and $F(r,\theta)=(r\cos\theta,r\sin\theta)$. The image is
the upper half-plane $V=\{y>0\}$, and the inverse is
\[ G(x,y)=\left(\sqrt{x^2+y^2},\ \arccos\frac{x}{\sqrt{x^2+y^2}}\right). \]
\end{example}

\begin{center}
\begin{tikzpicture}[scale=0.95]
  \begin{scope}
    \fill[plum, opacity=0.15] (0,0) rectangle (3,2.4);
    \draw[navy, ->, line width=0.6pt] (-0.3,0) -- (3.4,0) node[right, font=\small]{$r$};
    \draw[navy, ->, line width=0.6pt] (0,-0.3) -- (0,2.9) node[above, font=\small]{$\theta$};
    \draw[plum, dashed, line width=0.9pt] (0,2.4) -- (3,2.4);
    \node[font=\scriptsize, left] at (0,2.4) {$\pi$};
    \node[font=\small, text=plum] at (1.5,1.2) {$U$};
  \end{scope}
  \draw[-{Stealth[length=3mm]}, plum, line width=1.3pt] (3.8,1.6) to[bend left=20] node[above, font=\small]{$F$} (5.3,1.6);
  \draw[-{Stealth[length=3mm]}, coral, line width=1.3pt] (5.3,0.9) to[bend left=20] node[below, font=\small]{$G$} (3.8,0.9);
  \begin{scope}[xshift=8.2cm]
    \fill[plum, opacity=0.15] (-2.4,0) rectangle (2.4,2.4);
    \draw[navy, ->, line width=0.6pt] (-2.6,0) -- (2.7,0) node[right, font=\small]{$x$};
    \draw[navy, ->, line width=0.6pt] (0,-0.3) -- (0,2.9) node[above, font=\small]{$y$};
    \draw[plum, dashed, line width=0.9pt] (-2.4,0) -- (2.4,0);
    \node[font=\small, text=plum] at (1.2,1.4) {$V=\{y>0\}$};
  \end{scope}
\end{tikzpicture}
\end{center}

\begin{definition}{Local $C^1$-invertibility}
$F$ is \emph{locally $C^1$-invertible at $P$} if there is an open set $U_1\subset U$ containing $P$
such that $F$ is $C^1$-invertible on $U_1$.
\end{definition}

\subsection{The inverse mapping theorem}

\begin{theorem}[label=thm:inverse]{Inverse mapping theorem}
Let $F:U\to\R^n$ be a $C^1$-map and $P\in U$. If the Jacobian determinant $\Delta_F(P)\neq0$, then
$F$ is locally $C^1$-invertible at $P$.
\end{theorem}

\begin{intuition}
Near $P$, $F$ looks like its derivative, a linear map. A linear map is invertible exactly when its
determinant is nonzero. The theorem says that the nonlinear $F$ inherits this invertibility, at
least on a small neighbourhood. (The proof uses a fixed-point argument and is beyond this course.)
\end{intuition}

\begin{itemize}
  \item $\Delta_F(P)\neq0$ means the Jacobian matrix $J_F(P)$ is invertible.
  \item $J_F(P)$ represents $F'(P)$.
  \item So: \textbf{if $F'(P)$ is invertible, then $F$ is locally invertible at $P$.}
\end{itemize}

\begin{proposition}{Derivative of an inverse map}
If $F:U\to V$ has a $C^1$-inverse $G$, then for $Y=F(X)$,
\[ \boxed{\,G'(Y)=F'(X)^{-1}\,}\qquad\text{i.e.}\qquad J_G(Y)=J_F(X)^{-1}. \]
\end{proposition}

\begin{proof}
$G\circ F=\mathrm{id}$ and the identity map has derivative $I$. By the chain rule,
$I=(G\circ F)'(X)=G'(F(X))\circ F'(X)$. Similarly $F\circ G=\mathrm{id}$ gives $F'(X)\circ G'(Y)=I$.
So $G'(Y)$ is the inverse of $F'(X)$.
\end{proof}

\begin{example}{a map that is locally but not globally invertible}
Let $F(x,y)=\big(e^x\cos y,\;e^x\sin y\big)$. Then
\[ J_F(x,y)=\begin{pmatrix}e^x\cos y&-e^x\sin y\\ e^x\sin y&e^x\cos y\end{pmatrix}, \]
\[ \Delta_F(x,y)=e^{2x}(\cos^2y+\sin^2y)=e^{2x}\neq0 . \]
So $F$ is locally invertible at \emph{every} point. But $F(x,y+2\pi)=F(x,y)$: it is not injective,
so it has no global inverse on $\R^2$.

\textbf{Derivative of the local inverse at $F(0,\frac\pi2)=(0,1)$.}
$J_F(0,\frac\pi2)=\begin{pmatrix}0&-1\\1&0\end{pmatrix}$, so
$J_G(0,1)=\begin{pmatrix}0&-1\\1&0\end{pmatrix}^{-1}=\begin{pmatrix}0&1\\-1&0\end{pmatrix}$.
\end{example}

\begin{trap}
The slides write $\Delta_F(x,y)=e^x$ for this map. Computing the determinant gives
$e^x\cos y\cdot e^x\cos y+e^x\sin y\cdot e^x\sin y=e^{2x}$. The conclusion (never zero, so locally
invertible everywhere) is the same.
\end{trap}

\begin{reallife}{why ``locally'' matters}
In complex numbers this map is $z\mapsto e^z$. Every nonzero complex number has a logarithm, but
infinitely many of them (they differ by $2\pi i$). A calculator must choose one branch; this is the
same phenomenon as choosing an angle $\theta\in(-\pi,\pi]$ in polar coordinates.
\end{reallife}

\subsection{The implicit function theorem}

Consider a curve given by an equation $f(x,y)=c$, such as the circle $x^2+y^2=2$. We would like
to describe it, near a point, as a graph $y=\varphi(x)$.

\begin{definition}{Implicit function}
Suppose $f(a,b)=c$. If there is a differentiable function $y=\varphi(x)$, defined near $x=a$,
with $\varphi(a)=b$ and $f(x,\varphi(x))=c$, we say that $y=\varphi(x)$ is the function
\emph{determined implicitly} by $f$.
\end{definition}

\begin{example}{a map built for implicit functions}
Let $f:U\to\R$ be $C^1$ on an open $U\subset\R^2$, $(a,b)\in U$ with $D_2f(a,b)\neq0$. Then
$F(x,y)=(x,\,f(x,y))$ is locally invertible at $(a,b)$. Indeed
\[ J_F(x,y)=\begin{pmatrix}1&0\\ \pd fx&\pd fy\end{pmatrix},\qquad \Delta_F(a,b)=D_2f(a,b)\neq0, \]
and the inverse mapping theorem applies.
\end{example}

\begin{theorem}[label=thm:implicit]{Implicit function theorem}
Let $f$ be $C^1$ on an open set $U\subset\R^2$, $(a,b)\in U$, $f(a,b)=c$, and assume $D_2f(a,b)\neq0$.
Then there is an implicit function $y=\varphi(x)$, $C^1$ on some interval containing $a$, with
$\varphi(a)=b$ and $f(x,\varphi(x))=c$.
\end{theorem}

\begin{intuition}
$D_2f(a,b)\neq0$ means that, moving vertically through $(a,b)$, the value of $f$ changes. So for
each $x$ near $a$ there is exactly one nearby height $y$ where $f$ hits the value $c$ again. That
height is $\varphi(x)$. The map $F(x,y)=(x,f(x,y))$ packages this: it keeps $x$ and records the
value of $f$; inverting it means ``given $x$ and a value, find the $y$''.
\end{intuition}

\begin{proof}
\textbf{Step 1.} Define $F(x,y)=(x,f(x,y))$. By the previous example, $F$ has a local $C^1$-inverse
$G$ defined near $F(a,b)=(a,c)$.

\textbf{Step 2.} Since $F$ does not change the first coordinate, neither does $G$: write
$G(x,z)=(x,g(x,z))$ with $g$ of class $C^1$. Define $\varphi(x)=g(x,c)$.

\textbf{Step 3.} From $F(G(x,c))=(x,c)$ we get $F(x,\varphi(x))=(x,c)$, i.e. $f(x,\varphi(x))=c$.

\textbf{Step 4.} From $G(a,c)=(a,b)$ we get $\varphi(a)=b$. And $\varphi$ is $C^1$ because $g$ is.
\end{proof}

\begin{example}{the circle near $(1,1)$ and near $(-1,-1)$}
Let $f(x,y)=x^2+y^2$, $c=2$.
\begin{itemize}
  \item At $(1,1)$: $D_2f(1,1)=2\neq0$, so an implicit function exists; explicitly $y=\sqrt{2-x^2}$.
  \item At $(-1,-1)$: $D_2f(-1,-1)=-2\neq0$; explicitly $y=-\sqrt{2-x^2}$.
  \item At $(\sqrt2,0)$: $D_2f=0$. Indeed no function $y=\varphi(x)$ describes the circle there: for
  $x$ slightly less than $\sqrt2$ there are two values $\pm\sqrt{2-x^2}$, and for $x>\sqrt2$ none.
\end{itemize}
\end{example}

\begin{center}
\begin{tikzpicture}
\begin{axis}[bp, width=6.6cm, height=6.6cm, xmin=-2, xmax=2, ymin=-2, ymax=2,
   xtick={-1,1}, ytick={-1,1}, xlabel=$x$, ylabel=$y$, axis equal image]
  \addplot[plum!40, line width=1.1pt, domain=0:360, samples=90] ({1.4142*cos(x)},{1.4142*sin(x)});
  \addplot[coral, line width=2.2pt, domain=0.45:1.3, samples=30] {sqrt(2-x^2)};
  \addplot[cyanbp, line width=2.2pt, domain=-1.3:-0.45, samples=30] {-sqrt(2-x^2)};
  \node[dot] at (axis cs:1,1) {};
  \node[circle, fill=cyanbp, draw=white, inner sep=1.7pt] at (axis cs:-1,-1) {};
  \node[odot] at (axis cs:1.4142,0) {};
  \draw[navy, dashed] (axis cs:1.4142,-1.6) -- (axis cs:1.4142,1.6);
  \node[font=\scriptsize, text=coraldark, anchor=south west] at (axis cs:1.02,1.02) {$y=\sqrt{2-x^2}$};
  \node[font=\scriptsize, text=cyanbp, anchor=north east] at (axis cs:-1.02,-1.02) {$y=-\sqrt{2-x^2}$};
  \node[font=\scriptsize, text=navy, anchor=west, fill=white, inner sep=1pt] at (axis cs:1.45,0.35) {$D_2f=0$};
\end{axis}
\end{tikzpicture}

{\small\color{navy} Near $(1,1)$ and $(-1,-1)$ the circle is locally a graph (thick arcs). At
$(\sqrt2,0)$ the tangent is vertical and no local graph $y=\varphi(x)$ exists.}
\end{center}

\begin{keyidea}
The representation $y=\varphi(x)$ is \textbf{local}: it describes the curve only near the chosen
point, and it need not hold for all $x$.
\end{keyidea}

\subsection{Implicit differentiation}

\begin{extra}
Once we know $\varphi$ exists and is $C^1$, differentiate $f(x,\varphi(x))=c$ with the chain rule:
\[ D_1f(x,\varphi(x))+D_2f(x,\varphi(x))\,\varphi'(x)=0, \]
so
\[ \boxed{\,\varphi'(x)=-\frac{D_1f(x,\varphi(x))}{D_2f(x,\varphi(x))}\,} \]
The hypothesis $D_2f\neq0$ is exactly what allows the division.
\end{extra}

\begin{example}{implicit differentiation}
Let $f(x,y)=x^2y+3y^3x^4-4$ and $(a,b)=(1,1)$. Check: $f(1,1)=1+3-4=0$. Since
$D_2f=x^2+9y^2x^4$ gives $D_2f(1,1)=10\neq0$, an implicit function $y=\varphi(x)$ exists near $x=1$.
Differentiating $f(x,y)=0$ with respect to $x$ (treating $y$ as $\varphi(x)$):
\[ 2xy+x^2y'+12y^3x^3+9y^2y'x^4=0, \]
so
\[ \varphi'(x)=y'=-\frac{2xy+12y^3x^3}{x^2+9y^2x^4},\qquad \varphi'(1)=-\frac{2+12}{1+9}=-\frac75 . \]
\end{example}

\begin{example}{another implicit derivative}
Suppose $x^3+4y\sin(xy)=0$. Differentiating with respect to $x$ (product and chain rules):
\[ 3x^2+4y'\sin(xy)+4y\cos(xy)\,(y+xy')=0 . \]
Collecting the $y'$ terms:
\[ y'=-\frac{3x^2+4y^2\cos(xy)}{4\sin(xy)+4xy\cos(xy)} \]
whenever $4\sin(xy)+4xy\cos(xy)\neq0$.
\end{example}

\begin{reallife}{equations of state}
Thermodynamics is full of implicit functions. An equation of state $f(P,V,T)=0$ (for an ideal gas,
$PV-nRT=0$) links pressure, volume and temperature. Asking ``how does volume change with
temperature at constant pressure?'' is asking for an implicit derivative,
$\left(\pd VT\right)_P=-\dfrac{\partial f/\partial T}{\partial f/\partial V}=\dfrac{nR}{P}$.
\end{reallife}

\section{Summary}

\begin{keyidea}
\begin{itemize}
  \item $F$ is differentiable at $P$ if $F(P+H)=F(P)+LH+o(H)$ for a linear $L$; this $L$ is unique and is $F'(P)$.
  \item $F'(P)$ is represented by the Jacobian matrix $J_F(P)=(\partial f_i/\partial x_j)$.
  \item Differentials: $df=\sum f_{x_i}dx_i$; under $x=F(u)$, $dx=J_F\,du$ and
  $dV_x=\abs{\det J_F}\,dV_u$.
  \item Chain rule: $J_{G\circ F}(X)=J_G(F(X))\,J_F(X)$.
  \item Inverse mapping theorem: $\Delta_F(P)\neq0\Rightarrow F$ locally invertible, and $J_G=J_F^{-1}$.
  \item Implicit function theorem: $D_2f(a,b)\neq0\Rightarrow y=\varphi(x)$ near $a$, with $\varphi'=-D_1f/D_2f$.
\end{itemize}
\end{keyidea}

\section{Exam-style questions}

\begin{exercise}[1]
Compute the Jacobian matrix of $F(x,y)=\big(x^2y,\;x+y^3,\;\sin(xy)\big)$ and evaluate it at $(1,0)$.
\end{exercise}

\begin{exercise}[2]
Spherical coordinates are given by
\[ F(\rho,\varphi,\theta)=(\rho\sin\varphi\cos\theta,\;\rho\sin\varphi\sin\theta,\;\rho\cos\varphi). \]
Show that $\Delta_F=\rho^2\sin\varphi$, and explain what this says about the volume element.
\end{exercise}

\begin{exercise}[1]
Let $F(x,y)=(x+y,\;x-y)$ and $G(u,v)=(uv,\;u^2-v^2)$. Compute $J_{G\circ F}(1,1)$ with the chain rule,
then check by composing first.
\end{exercise}

\begin{exercise}[2]
Let $z=x^2y$ with $x=r\cos\theta$, $y=r\sin\theta$. Express $dz$ in terms of $dr$ and $d\theta$ using
differentials, and check against $z$ written in polar coordinates.
\end{exercise}

\begin{exercise}[2]
Let $F(x,y)=(x+y^2,\;y+x^2)$.
\begin{enumerate}[label=(\alph*)]
  \item Find all points where the inverse mapping theorem guarantees local invertibility.
  \item Compute the derivative of the local inverse at $F(1,1)=(2,2)$.
\end{enumerate}
\end{exercise}

\begin{exercise}[2]
The folium of Descartes is $x^3+y^3=6xy$.
\begin{enumerate}[label=(\alph*)]
  \item Show that near $(3,3)$ the curve is the graph of a $C^1$ function $y=\varphi(x)$, and find
  the tangent line there.
  \item Find the points of the curve where the implicit function theorem (solving for $y$) does not apply.
\end{enumerate}
\end{exercise}

\begin{exercise}[3]
For the polar map $F(r,\theta)=(r\cos\theta,r\sin\theta)$, use the formula $J_G=J_F^{-1}$ at
$(r,\theta)=(2,\frac\pi2)$ to find $\pd rx,\pd ry,\pd\theta x,\pd\theta y$ at the point $(x,y)=(0,2)$. Check
two of them by differentiating $r=\sqrt{x^2+y^2}$ and $\theta=\arccos(x/r)$ directly.
\end{exercise}

\begin{exercise}[3]
\begin{enumerate}[label=(\alph*)]
  \item Let $A$ be an $m\times n$ matrix, $B\in\R^m$ and $F(X)=AX+B$. Show directly from the definition
  that $F'(P)=A$ for every $P$.
  \item Let $f(X)=X\cdot X$ on $\R^n$. Show directly from the definition that $f'(P)H=2P\cdot H$.
\end{enumerate}
\end{exercise}
